% PTE Core 模考 3B —— 阅读 错题与详解·整理卷  版本 000（自包含）
% 编译: xelatex 本文件.tex   （需 Noto CJK + TeX Gyre Heros 字体）
% 本版本无外部图片，单文件即可复现。
\documentclass[11pt]{article}

% --------- 样式（原 ptestyle.sty，已内联）---------
% ============================================================
%  ptestyle.sty  —  PTE 模考错题整理卷 共享样式
%  编译引擎: XeLaTeX  (中英混排, 需要 Noto CJK 字体)
% ============================================================

\usepackage[a4paper,margin=2cm,headsep=4mm]{geometry}
\usepackage{fontspec}
\usepackage{xeCJK}
\usepackage[table]{xcolor}
\usepackage[normalem]{ulem}        % \sout 删除线
\usepackage{tabularx}
\usepackage{array}
\usepackage{ragged2e}
\usepackage{enumitem}
\usepackage[most]{tcolorbox}
\usepackage{fancyhdr}
\usepackage{graphicx}
\usepackage{pifont}                % \ding{51}=✓ \ding{55}=✗（阅读对错标注）
\usepackage{needspace}
\usepackage{tikz}

% ---------- 字体 ----------
\setmainfont{TeX Gyre Heros}            % 拉丁: Helvetica 风格(纤细清晰)
\setCJKmainfont{Noto Sans CJK SC}
\setCJKsansfont{Noto Sans CJK SC}
\newfontfamily\cjkserif{Noto Serif CJK SC}
\linespread{1.18}
\setlength{\parindent}{0pt}
\setlength{\parskip}{4pt}

% ---------- 配色 ----------
\definecolor{cWrong}{HTML}{D32F2F}     % 红  = 修改 / 错误
\definecolor{cFix}{HTML}{1B7F3B}       % 绿  = 正确建议
\definecolor{cDelete}{HTML}{1565C0}    % 蓝  = 删除
\definecolor{cInsert}{HTML}{7B1FA2}    % 紫  = 插入
\definecolor{cPartial}{HTML}{E07B00}   % 橙  = 部分得分
\definecolor{cNote}{HTML}{0277BD}      % 蓝  = 注解
\definecolor{cAccent}{HTML}{16A085}    % 主题青绿(平台主色)
\definecolor{cWrongBg}{HTML}{FBE0E0}
\definecolor{cDelBg}{HTML}{DCE8FB}
\definecolor{cInsBg}{HTML}{EEE0F5}
\definecolor{cHead}{HTML}{20303A}
\definecolor{cGrayBox}{HTML}{F4F6F7}
\definecolor{cModelBg}{HTML}{EAF7EF}
\definecolor{cYouBg}{HTML}{FFF7F5}
\definecolor{cNoteBg}{HTML}{EAF3FA}
\definecolor{cRule}{HTML}{D7DEE2}

% ---------- 行内批改宏 (复刻平台标注) ----------
% \fix{原词}{正确}  红色删除线 + 绿色上标改正  (修改/拼写/空格)
\newcommand{\fix}[2]{%
  \colorbox{cWrongBg}{\textcolor{cWrong}{\sout{#1}}}%
  \,\textsuperscript{\textcolor{cFix}{\footnotesize #2}}}
% \del{原词}  蓝色删除线  (应删除)
\newcommand{\del}[1]{%
  \colorbox{cDelBg}{\textcolor{cDelete}{\sout{#1}}}%
  \,\textsuperscript{\textcolor{cDelete}{\footnotesize 删}}}
% \ins{词}  紫色下划线  (应插入)
\newcommand{\ins}[1]{%
  \textsuperscript{\textcolor{cInsert}{\footnotesize $\wedge$}}%
  \colorbox{cInsBg}{\textcolor{cInsert}{#1}}}
% \miss{词}  灰色删除线  (口语漏读 / 未作答)
\newcommand{\miss}[1]{\textcolor{gray}{\sout{#1}}}
% \add{原词}{改正}  橙色波浪线 + "补"标签  (平台漏标、本卷补标的错误)
\definecolor{cAdd}{HTML}{E65100}
\newcommand{\add}[2]{%
  \textcolor{cAdd}{\uwave{#1}}%
  \,\textsuperscript{\textcolor{cAdd}{\footnotesize 补：#2}}}
% 高亮(口语发音): 好/中/差
\newcommand{\spk}[2]{\textcolor{#1}{#2}}

% ---------- 评分单元格着色 ----------
\newcommand{\sfull}[1]{\textcolor{cFix}{\bfseries #1}}     % 满分
\newcommand{\spart}[1]{\textcolor{cPartial}{\bfseries #1}} % 部分
\newcommand{\szero}[1]{\textcolor{cWrong}{\bfseries #1}}   % 零/低分

% ---------- 分数小标签 ----------
\newtcbox{\chip}[1][cAccent]{on line, boxsep=2pt, left=4pt, right=4pt, top=1pt,
  bottom=1pt, arc=3pt, colback=#1, colframe=#1, coltext=white, nobeforeafter,
  fontupper=\bfseries\footnotesize}

% ---------- 题块小标题 ----------
\newcommand{\ptesub}[1]{%
  \par\needspace{2\baselineskip}\smallskip
  {\color{cAccent}\rule[-0.2ex]{3pt}{1.05em}}\hspace{6pt}%
  {\bfseries\color{cHead}#1}\par\nopagebreak\smallskip}

% ---------- 题块外框 ----------
% \begin{qblock}{标号·题型}{右上信息(得分/用时)} ... \end{qblock}
\newtcolorbox{qblock}[2]{
  breakable, enhanced, arc=2.5pt,
  colback=white, colframe=cRule, boxrule=0.6pt,
  left=11pt, right=11pt, top=8pt, bottom=10pt,
  toptitle=3pt, bottomtitle=3pt, lefttitle=11pt, righttitle=11pt,
  colbacktitle=cHead, coltitle=white, fonttitle=\bfseries,
  title={#1\hfill\normalfont\bfseries\small #2}}

% ---------- 文本块: 题干 / 范文 / 你的作答 / 建议 ----------
\newtcolorbox{promptbox}{enhanced, breakable, colback=cGrayBox, colframe=cGrayBox,
  arc=2pt, left=8pt, right=8pt, top=6pt, bottom=6pt, boxrule=0pt}
\newtcolorbox{modelbox}{enhanced, breakable, colback=cModelBg, colframe=cModelBg,
  borderline west={3pt}{0pt}{cFix}, arc=1pt, left=8pt, right=8pt, top=6pt, bottom=6pt, boxrule=0pt}
\newtcolorbox{youbox}{enhanced, breakable, colback=cYouBg, colframe=cYouBg,
  borderline west={3pt}{0pt}{cWrong}, arc=1pt, left=8pt, right=8pt, top=6pt, bottom=6pt, boxrule=0pt}
\newtcolorbox{notebox}{enhanced, breakable, colback=cNoteBg, colframe=cNoteBg,
  borderline west={3pt}{0pt}{cNote}, arc=1pt, left=8pt, right=8pt, top=6pt, bottom=6pt, boxrule=0pt}

% ---------- 评分表 ----------
% 用法见正文; 这里给表头宏
\newcommand{\scorehead}{%
  \rowcolors{2}{cGrayBox}{white}
  \arrayrulecolor{cRule}}

% ---------- 章节大标题 ----------
\newcommand{\ptesection}[3]{% #1=中文名 #2=英文名 #3=该项分数
  \addcontentsline{toc}{section}{#1}%
  \needspace{6\baselineskip}
  \begin{tcolorbox}[enhanced, colback=cAccent, colframe=cAccent, sharp corners,
     left=14pt, right=14pt, top=10pt, bottom=10pt, boxrule=0pt]
    {\fontsize{20}{24}\selectfont\bfseries\color{white} #1\ \ }%
    {\large\color{white!85} #2}\hfill
    {\large\color{white}本项得分\ \ }{\fontsize{22}{24}\selectfont\bfseries\color{white} #3}%
    \ {\color{white!85}/ 90}
  \end{tcolorbox}}

% ---------- 题型小节分隔 (口语 RA/RS/DI/RTS/ASQ 等) ----------
\newcommand{\ptetask}[2]{%
  \addcontentsline{toc}{subsection}{#1\ #2}%
  \needspace{4\baselineskip}\medskip
  \begin{tcolorbox}[enhanced, colback=cHead, colframe=cHead, sharp corners,
     left=12pt, right=12pt, top=5pt, bottom=5pt, boxrule=0pt]
    {\large\bfseries\color{white} #1}\ \ {\color{white!80} #2}
  \end{tcolorbox}\smallskip}

% ---------- 口语 AI 语音识别着色 (复刻平台: 优/良/差 + 停顿) ----------
\definecolor{cSpkGood}{HTML}{2E9E5B}   % 优 绿
\definecolor{cSpkOk}{HTML}{C77F00}     % 良 黄/琥珀
\definecolor{cSpkBad}{HTML}{E0352B}    % 差 红
\newcommand{\sg}[1]{\textcolor{cSpkGood}{#1}}   % 优
\newcommand{\sy}[1]{\textcolor{cSpkOk}{#1}}     % 良
\newcommand{\sr}[1]{\textcolor{cSpkBad}{#1}}    % 差
\newcommand{\smiss}[1]{\textcolor{cSpkBad}{\sout{#1}}}  % 漏读 (红删除线)
\newcommand{\pse}{{\color{gray}\,/\,}}          % 停顿
\newcommand{\asrlegend}{{\footnotesize\color{gray}AI 语音识别\quad
  {\color{cSpkGood}\textbullet}\,优\ \ {\color{cSpkOk}\textbullet}\,良\ \
  {\color{cSpkBad}\textbullet}\,差\ \ {\color{gray}/}\,停顿\ \
  {\color{cSpkBad}\sout{词}}\,漏读}}
\newtcolorbox{asrbox}{enhanced, breakable, colback=white, colframe=cRule,
  boxrule=0.6pt, arc=1pt, left=8pt, right=8pt, top=6pt, bottom=6pt}

% ---------- 中文翻译行 ----------
\newcommand{\zh}[1]{{\footnotesize\textcolor{cNote}{\textbf{译}}\ \textcolor{gray}{#1}}}

% ---------- 阅读填空 / 选择 对错标注 ----------
\newcommand{\rok}[1]{{\color{cFix}\ding{51}\,#1}}                 % 选对(绿勾 + 词)
\newcommand{\rno}[2]{{\color{cWrong}\ding{55}\,\sout{#1}}\,{\footnotesize\color{cFix}(答案:\,#2)}}  % 选错: 你的(红删) + (答案:正确)
\newcommand{\rblank}[1]{{\color{cFix}\underline{\,#1\,}}}         % 直接给空的正确答案(无你的作答时)

% ---------- 听力专用 ----------
\newcommand{\hiw}[2]{{\color{cWrong}\uline{#1}}\,{\footnotesize\color{cFix}(听:#2)}} % HIW: 屏幕错词(红下划线)+ 实际读音
\newcommand{\lhit}[1]{{\color{cFix}#1}}                           % WFD: 命中的词(绿)
\newcommand{\lmiss}[1]{{\color{cWrong}\sout{#1}}}                 % WFD: 多余/错误的词(红删, 1B 配色)
% ---- WFD 逐词批改 · 本套 2B 平台配色（与 1B 相反）：绿=命中 / 红括号=漏写 / 灰删除线=多写或听错 ----
% 说明: 2B 平台把"漏写"标红(加括号)、"多写/听错"标灰删除线; 本卷忠于原图按此复刻。
\newcommand{\womit}[1]{{\color{cWrong}(#1)}}                      % WFD: 漏写(正确答案有、你没写)——红色括号
\newcommand{\wbad}[1]{{\color{gray}\sout{#1}}}                    % WFD: 多写/听错(你写了但错)——灰删除线
\newcommand{\wfdlegend}{{\footnotesize\color{gray}WFD 配色：\
  {\color{cFix}命中}\ \ {\color{cWrong}(漏写)}\ \ {\color{gray}\sout{多写/听错}}}}

% ---------- 颜色图例 ----------
\newcommand{\ptelegend}{%
  \begin{tcolorbox}[enhanced, colback=cGrayBox, colframe=cRule, boxrule=0.5pt,
    arc=2pt, left=10pt, right=10pt, top=6pt, bottom=6pt]
  {\bfseries\color{cHead}颜色说明\quad}
  \fix{改}{正}~修改/拼写\quad
  \del{删}~应删除\quad
  \ins{插}~应插入\quad
  \add{补}{改正}~平台漏标(本卷补标)\quad
  {\color{cFix}\rule{1em}{0.6em}}~范文/正确\quad
  {\color{cNote}\rule{1em}{0.6em}}~AI 建议
  \end{tcolorbox}}

% ---------- 页眉页脚 ----------
\pagestyle{fancy}
\fancyhf{}
\renewcommand{\headrulewidth}{0.4pt}
\renewcommand{\footrulewidth}{0pt}
\fancyhead[L]{\footnotesize\color{cHead}\bfseries PTE Core 模考 3B}
\fancyhead[R]{\footnotesize\color{gray} 错题与详解整理卷}
\fancyfoot[C]{\footnotesize\color{gray}\thepage}

\begin{document}

% --------- 封面（原 cover.tex）---------
% ---------- 封面 / 成绩总览 ----------
\definecolor{mListen}{HTML}{1F3A5F}
\definecolor{mRead}{HTML}{C2D70B}
\definecolor{mSpeak}{HTML}{4D4D4D}
\definecolor{mWrite}{HTML}{A01B7D}

\begin{titlepage}
\thispagestyle{empty}
\centering
\vspace*{1.4cm}

{\fontsize{30}{36}\selectfont\bfseries\color{cHead} PTE Core 模考 3B}\\[4mm]
{\fontsize{17}{20}\selectfont\color{cAccent}\bfseries 错题与详解 · 整理卷}\\[3mm]
{\color{gray} APEUni / 猩际 PTE 模考　·　提交时间 2026-07-05 13:36}\\[1.2cm]

% ---- 总分 ----
\begin{tcolorbox}[width=0.62\textwidth, halign=center, enhanced,
  colback=cAccent, colframe=cAccent, arc=6pt, boxrule=0pt, top=10pt, bottom=10pt]
  {\color{white}\large 总分 Overall}\\[2pt]
  {\fontsize{46}{50}\selectfont\bfseries\color{white} 53}
\end{tcolorbox}

\vspace{1.0cm}

% ---- 四项分数条形图 ----
\begin{tcolorbox}[width=0.84\textwidth, enhanced, colback=white, colframe=cRule,
  boxrule=0.8pt, arc=4pt, left=18pt, right=18pt, top=14pt, bottom=14pt]
\setlength{\tabcolsep}{6pt}\renewcommand{\arraystretch}{1.7}
\sffamily
\begin{tabular}{@{}r@{\hskip 8pt}l@{\hskip 10pt}l@{}}
 \bfseries Listening 听力 & {\color{mListen}\rule{6.11cm}{0.42cm}} & {\bfseries\large 55}\\
 \bfseries Reading 阅读   & {\color{mRead}\rule{4.67cm}{0.42cm}}   & {\bfseries\large 42}\\
 \bfseries Speaking 口语  & {\color{mSpeak}\rule{4.67cm}{0.42cm}}  & {\bfseries\large 42}\\
 \bfseries Writing 写作   & {\color{mWrite}\rule{8.11cm}{0.42cm}}  & {\bfseries\large 73}\\
\end{tabular}\\[2mm]
{\footnotesize\color{gray} 满分 90　|　刻度: 条长 = 得分 / 90}
\end{tcolorbox}

\vfill
\begin{flushleft}\small\color{gray}
\textbf{说明:}本卷由本次模考的题目与"详解"截图逐题誊写、重新排版而成,完整保留每道题的
\emph{题干 / 原文、范文、你的作答、逐项评分、彩色批改与 AI 中文建议}。\\
所有错误按平台标注用颜色标出(见下方图例)。原始"详解"PDF 不可复制、仅截图,本卷内容均为人工对照誊录,若与原图有出入以原图为准。
\end{flushleft}
\vspace{4mm}
\ptelegend
\end{titlepage}

% --------- 阅读正文（原 sec_reading_body.tex）---------
% ===================== 阅读 Reading =====================
\ptesection{阅读}{Reading · FIB + MCM-R + RO + FIBD\&D + MCS-R}{42}

\vspace{1mm}
{\small\color{gray} 本部分共 \textbf{16 题}：\textbf{FIB 读写填空}（下拉选词）Q1--5、\textbf{MCM-R 多选}Q6--7、\textbf{RO 段落重排}Q8--9、\textbf{FIBD\&D 读填空}（拖词入空）Q10--14、\textbf{MCS-R 单选}Q15--16。
每题保留文章原文、你的作答（空处对错与平台一致：\rok{绿勾}=选对，{\color{cWrong}\ding{55}}\,红叉=选错并给正确答案）、选项/词库（正确项加粗）、你的答案与得分；并为每篇文章附\textbf{中文翻译}。本套阅读总分 \textbf{42 / 90}。}
\vspace{1mm}

{\small\color{gray}\textbf{说明：}本套复盘截图里每个空后面都有一个「？」小图标（点开显示解析气泡），但截图本身只在个别题上恰好停留在展开状态——\textbf{仅 Q13 的解析被完整截了下来}，其余题目的解析气泡未展开，故从略（如需可到平台点「？」或「查看原题」补看）。}
\vspace{2mm}

\newcolumntype{Z}{>{\RaggedRight\arraybackslash}X}

% ---------------- Q1 ----------------
\begin{qblock}{Q1\quad FIB \#384\ \ Starvation}{读写填空 · 评分 \textcolor{cAccent}{\bfseries 3 / 5}}
\ptesub{文章 Passage（含填空作答）}
\begin{promptbox}
Over weeks and months, malnutrition can result in specific diseases, like anemia when people don't get enough iron or beriberi if they don't get adequate thiamine. A \rno{distinctive}{severe} lack of food for a prolonged period --- not enough calories of any sort to keep up with the body's energy needs --- is starvation. The body's reserve resources are \rok{depleted}. The result is substantial weight loss, wasting away of the body's tissues and eventually death. When faced with starvation, the body fights back. The first day without food is a lot like the overnight fast between dinner one night and breakfast the next morning. Energy levels are low but \rno{feed}{pick} up with a morning meal. Within days, faced with nothing to eat, the body begins feeding on itself. Metabolism slows; the body cannot regulate its temperature; kidney function is impaired and the immune system \rok{weakens}. When the body uses its reserves to provide basic energy needs, it can no longer supply necessary nutrients to vital organs and tissues. The heart, lungs, ovaries and testes shrink. Muscles shrink and people feel weak. Body temperature drops and people can feel chilled. People can become \rok{irritable}, and they become difficult to concentrate.
\end{promptbox}
\zh{经过数周乃至数月，营养不良会导致特定疾病，比如缺铁引起的贫血，或缺乏足量硫胺素引起的脚气病。长期严重缺乏食物——任何种类的热量都不足以维持身体能量需求——就是饥饿。此时身体的储备资源被耗尽，结果是体重大幅下降、身体组织逐渐消耗，最终死亡。面对饥饿，身体会做出反抗：不进食的第一天很像晚饭到次日早饭之间的整夜禁食，能量水平较低，但早餐后又能恢复。几天之内，面对无物可食，身体开始"吃自己"：代谢减慢，身体无法调节体温，肾功能受损，免疫系统减弱。当身体动用储备来满足基本能量需求时，就无法再向重要器官和组织供应必需的营养——心脏、肺、卵巢和睾丸都会萎缩，肌肉萎缩、人变得虚弱，体温下降、人会感到发冷。人还可能变得易怒，难以集中注意力。}
\ptesub{选项 Choices（每空 4 选 1，正确项加粗）}
{\small
1.\ \textbf{severe},\ distinguishing,\ proper,\ distinctive\\
2.\ obsoleted,\ \textbf{depleted},\ pelleted,\ deleted\\
3.\ feed,\ come,\ chill,\ \textbf{pick}\\
4.\ deepens,\ deafens,\ \textbf{weakens},\ surpasses\\
5.\ \textbf{irritable},\ commutable,\ indisputable,\ transportable}
\par\smallskip
\textbf{你的答案：}\ {\color{cWrong}distinctive}, depleted, {\color{cWrong}feed}, weakens, irritable\qquad\textbf{本题得分：}\ {\color{cAccent}\bfseries 3 / 5}\quad{\footnotesize\color{gray}（用时 02:28；第 1 空 distinctive→severe、第 3 空 feed→pick）}
\end{qblock}

% ---------------- Q2 ----------------
\begin{qblock}{Q2\quad FIB \#380\ \ Death Sentence}{读写填空 · 评分 \textcolor{cAccent}{\bfseries 2 / 5}}
\ptesub{文章 Passage（含填空作答）}
\begin{promptbox}
You are more likely to be \rok{sentenced} to death if you are a member of a minority group within a state that executes defendants. The death penalty disproportionately affects members of racial, ethnic and religious minorities, as well as those living in poverty. In the US, there's extensive evidence of racial \rno{appearance}{bias} on death row. The race of the victim remains the single most reliable factor in \rok{determining} whether a defendant will be given a death sentence. African American defendants are three times more likely to receive the death penalty than white defendants, where the victims are white. Serious mental health issues are also common in defendants sent to death row. At least one in ten prisoners \rno{persecuted}{executed} in the US between 1977 and 2007 had experienced severe mental health problems that meant they were literally unable to comprehend the crime they were \rno{persuaded}{alleged} to have committed, and unable to understand the terms of their sentence and imminent execution.
\end{promptbox}
\zh{如果你是执行死刑的州里少数群体的一员，你被判死刑的可能性更高。死刑不成比例地影响着种族、族裔和宗教上的少数群体，以及生活在贫困中的人。在美国，有大量证据表明死囚群体中存在种族偏见。受害者的种族仍然是决定被告是否会被判死刑的最可靠单一因素——当受害者是白人时，非裔美国被告被判死刑的可能性是白人被告的三倍。严重的心理健康问题在死囚犯中也很常见：1977年至2007年间，美国至少有十分之一被处决的囚犯曾经历严重的心理健康问题，以至于他们实际上根本无法理解自己被指控所犯下的罪行，也无法理解判决的内容和即将到来的处决。}
\ptesub{选项 Choices（每空 4 选 1，正确项加粗）}
{\small
1.\ penalized,\ blamed,\ complained,\ \textbf{sentenced}\\
2.\ \textbf{bias},\ equality,\ appearance,\ background\\
3.\ \textbf{determining},\ adjoining,\ undermining,\ examining\\
4.\ electrocuted,\ persecuted,\ \textbf{executed},\ captured\\
5.\ \textbf{alleged},\ acclaimed,\ persuaded,\ claimed}
\par\smallskip
\textbf{你的答案：}\ sentenced, {\color{cWrong}appearance}, determining, {\color{cWrong}persecuted, persuaded}\qquad\textbf{本题得分：}\ {\color{cAccent}\bfseries 2 / 5}\quad{\footnotesize\color{gray}（用时 02:13；第 2 空 appearance→bias、第 4 空 persecuted→executed、第 5 空 persuaded→alleged）}
\end{qblock}

% ---------------- Q3 ----------------
\begin{qblock}{Q3\quad FIB \#381\ \ Deception}{读写填空 · 评分 \textcolor{cAccent}{\bfseries 3 / 5}}
\ptesub{文章 Passage（含填空作答）}
\begin{promptbox}
Deception refers to the act of \rok{encouraging} people to believe information that is not true. Lying is a common form of deception, stating something known to be untrue with the intent to deceive. While most people are generally honest, even those who \rok{subscribe} to honesty engage in deception sometimes. Studies show that the average person lies several times a day. Some of those lies are big (`I've never cheated on you!') but more often, they are little white lies (`That dress looks fine') deployed to \rok{avoid} uncomfortable situations or spare someone's feelings. Trust is the bedrock of social life at all levels, from romance and parenting to national government. Deception always \rno{undertakes}{undermines} it. Because truth is so essential to the human enterprise, which relies on a shared view of reality, the default assumption most people have is that others are truthful in their communications and dealings. Most cultures have powerful social \rno{ejections}{sanctions} against lying.
\end{promptbox}
\zh{欺骗是指让人们相信不真实信息的行为。撒谎是欺骗的一种常见形式，即明知某事不实、却出于欺骗意图而陈述它。虽然大多数人总体上是诚实的，但即便是信奉诚实的人有时也会进行欺骗。研究显示，普通人平均每天会撒谎好几次。其中有些谎言很大("我从没有出轨过！")，但更多时候是善意的小谎言("这条裙子看起来挺好")，用来避免尴尬处境或顾及他人感受。信任是社会生活各个层面的基石，从恋爱、为人父母到国家治理概莫能外。欺骗总是会破坏信任。因为真相对依赖共享现实观的人类社会而言至关重要，大多数人默认的假设是：他人在交流与交往中是诚实的。大多数文化都有强有力的社会性制裁来对抗撒谎。}
\ptesub{选项 Choices（每空 4 选 1，正确项加粗）}
{\small
1.\ discouraging,\ forbidding,\ detecting,\ \textbf{encouraging}\\
2.\ describe,\ prescribe,\ inscribe,\ \textbf{subscribe}\\
3.\ contest,\ illuminate,\ disguise,\ \textbf{avoid}\\
4.\ \textbf{undermines},\ underscores,\ undertakes,\ underwrites\\
5.\ ejections,\ \textbf{sanctions},\ fractions,\ inductions}
\par\smallskip
\textbf{你的答案：}\ encouraging, subscribe, avoid, {\color{cWrong}undertakes, ejections}\qquad\textbf{本题得分：}\ {\color{cAccent}\bfseries 3 / 5}\quad{\footnotesize\color{gray}（用时 02:06；第 4 空 undertakes→undermines、第 5 空 ejections→sanctions）}
\end{qblock}

% ---------------- Q4 ----------------
\begin{qblock}{Q4\quad FIB \#373\ \ Information Revolution}{读写填空 · 评分 \textcolor{cAccent}{\bfseries 1 / 5}}
\ptesub{文章 Passage（含填空作答）}
\begin{promptbox}
Some have begun to call it the Information Revolution. Technological changes brought dramatic new \rok{options} to Americans living in the 1990s. From the beginning of the decade until the end, new forms of entertainment, commerce, research, work, and communication became commonplace in the United States. The \rno{unremitting}{driving} force behind much of this change was an innovation popularly known as the Internet. Personal computers had become widespread by the end of the 1980s. Through a device called a modem, individual users could link their computer to a \rno{volume}{wealth} of information using conventional phone lines. What lay beyond the individual computer was a vast domain of information known as cyberspace. Upon its \rno{relief}{release} in 1983 the Apple ``Lisa'' computer was supposed to revolutionize personal computing. But interest in ``Lisa'' was minimal due to its nearly \$10,000 price tag and the introduction of the much more \rno{advanced}{affordable} ``Macintosh'' a year later.
\end{promptbox}
\zh{有人已经开始称之为"信息革命"。技术变革为20世纪90年代的美国人带来了戏剧性的全新选择。从这十年之初到末尾，新形式的娱乐、商业、研究、工作和通讯在美国变得司空见惯。推动这一变化的主要力量，是一项被普遍称为"互联网"的创新。到20世纪80年代末，个人电脑已经普及。通过一种叫做调制解调器的设备，个人用户可以借助传统电话线将自己的电脑连接到海量的信息资源上。个人电脑之外，还存在一个被称为"网络空间"的广阔信息领域。1983年发布的苹果"Lisa"电脑本应彻底革新个人计算，但由于其近一万美元的高价，加上一年后推出的价格更加亲民的"Macintosh"，人们对"Lisa"的兴趣微乎其微。}
\ptesub{选项 Choices（每空 4 选 1，正确项加粗）}
{\small
1.\ challenges,\ puzzles,\ \textbf{options},\ confusion\\
2.\ unremitting,\ uninspiring,\ \textbf{driving},\ insinuating\\
3.\ magnitude,\ bulk,\ \textbf{wealth},\ volume\\
4.\ relief,\ \textbf{release},\ publication,\ emission\\
5.\ convenient,\ \textbf{affordable},\ advanced,\ formidable}
\par\smallskip
\textbf{你的答案：}\ options, {\color{cWrong}unremitting, volume, relief, advanced}\qquad\textbf{本题得分：}\ {\color{cAccent}\bfseries 1 / 5}\quad{\footnotesize\color{gray}（用时 02:53；仅第 1 空 options 正确，其余四空全错）}
\end{qblock}

% ---------------- Q5 ----------------
\begin{qblock}{Q5\quad FIB \#372\ \ Dentistry}{读写填空 · 评分 \textcolor{cAccent}{\bfseries 2 / 4}}
\ptesub{文章 Passage（含填空作答）}
\begin{promptbox}
Dentistry is a profession \rno{dealt}{concerned} with the prevention and treatment of oral disease. Dentistry also encompasses the treatment and correction of malformation of the jaws, misalignment of the teeth, and birth anomalies of the oral cavity such as cleft palate. Dentistry, in some form, has been \rok{practiced} since ancient times. For example, Egyptian skulls dating from 2900 to 2750 BCE contain evidence of small holes in the jaw in the vicinity of a tooth's roots. Such holes are believed to have been \rno{laminated}{drilled} to drain abscesses. In addition, accounts of dental treatment appear in Egyptian scrolls dating from 1500 BCE. It is thought that the Egyptians practiced oral surgery perhaps as early as 2500 BCE, although evidence for this is minimal. An early attempt at tooth \rok{replacement} dates to Phoenicia (modern Lebanon) around 600 BCE, where missing teeth were replaced with animal teeth and were bound into place with cord.
\end{promptbox}
\zh{牙科是一门致力于预防和治疗口腔疾病的专业。牙科还涵盖了对颌骨畸形、牙齿错位以及诸如唇腭裂等口腔先天异常的治疗与矫正。某种形式的牙科诊疗自古以来就已存在：例如，公元前2900年至2750年的埃及头骨上有证据显示，牙根附近的颌骨上有小孔，人们认为这些小孔是被钻开以引流脓肿的。此外，公元前1500年的埃及卷轴中也记载了牙科治疗的内容。据推测，埃及人也许早在公元前2500年就已经进行口腔手术，尽管这方面的证据很少。最早尝试补牙的记录可追溯到约公元前600年的腓尼基（今黎巴嫩），当时缺失的牙齿会用动物牙齿替代，并用绳子绑扎固定。}
\ptesub{选项 Choices（每空 4 选 1，正确项加粗）}
{\small
1.\ agreed,\ dealt,\ \textbf{concerned},\ taken\\
2.\ criticized,\ replaced,\ \textbf{practiced},\ abandoned\\
3.\ fluctuated,\ laminated,\ \textbf{drilled},\ sealed\\
4.\ reparation,\ sacrament,\ restitution,\ \textbf{replacement}}
\par\smallskip
\textbf{你的答案：}\ {\color{cWrong}dealt}, practiced, {\color{cWrong}laminated}, replacement\qquad\textbf{本题得分：}\ {\color{cAccent}\bfseries 2 / 4}\quad{\footnotesize\color{gray}（用时 01:38；第 1 空 dealt→concerned、第 3 空 laminated→drilled）}
\end{qblock}

% ---------------- Q6 ----------------
\begin{qblock}{Q6\quad MCM-R \#118\ \ Taste Sensitivity}{多选 · 评分 \textcolor{cAccent}{\bfseries 1 / 2}}
\ptesub{文章 Passage}
\begin{promptbox}
What you eat influences your taste for what you might want to eat next. So claims a University of California, Riverside, study performed on fruit flies. The study offers a better understanding of neurophysiological plasticity of the taste system in flies. To maintain ideal health, animals require a balanced diet with optimum amounts of different nutrients. Macronutrients like carbohydrates and proteins are essential; indeed, an unbalanced intake of these nutrients can be detrimental to health. Flies require macronutrients such as sugars and amino acids for survival. They use the gustatory system, the sensory system responsible for the perception of taste, to sense these nutrients and begin feeding.

In their experiments in the lab, the researchers Anindya Ganguly and Manali Dey, led by Anupama Dahanukar, fed adult flies different diets: a balanced diet, a sugar-reduced and protein-enriched diet, and a sugar-enriched and protein-depleted diet. They ensured that all three diets were similar in total calorie content and tested the flies daily for a week to examine modifications in their food choice and taste sensitivity. The researchers report that diet affects dopamine and insulin signaling in the brain, which, in turn, affects the flies' peripheral sensory response, which is comprised of neurons directly involved in detecting external stimuli. This response then influences what the flies eat next.
\end{promptbox}
\zh{你所吃的食物会影响你接下来想吃什么口味的东西——加州大学河滨分校一项针对果蝇的研究如是说。该研究让我们更好地理解了果蝇味觉系统的神经生理可塑性。为保持理想的健康状态，动物需要摄入均衡的饮食，各种营养素的量都要恰到好处。碳水化合物和蛋白质这类宏量营养素是必需的；事实上，这些营养素摄入不均衡会损害健康。果蝇的生存需要糖和氨基酸等宏量营养素，它们利用负责感知味道的感觉系统——味觉系统——来感知这些营养素并开始进食。在实验室实验中，由 Anupama Dahanukar 带领的研究人员 Anindya Ganguly 和 Manali Dey 给成年果蝇喂食了三种不同的饮食：均衡饮食、低糖高蛋白饮食，以及高糖低蛋白饮食。他们确保这三种饮食的总热量相近，并每天对果蝇进行为期一周的测试，以观察它们食物选择和味觉敏感度的变化。研究人员报告称，饮食会影响大脑中多巴胺和胰岛素的信号传导，进而影响果蝇的外周感觉反应——该反应由直接参与探测外部刺激的神经元构成。这种反应接下来又会影响果蝇下一步吃什么。}
\ptesub{题目 Question}
\begin{notebox}Which of the following statements about the study are true?\end{notebox}
\ptesub{选项 Choices（正确项绿色加粗）}
{\small
A) What you eat has little to do with what you want to eat next.\\
B) An unbalanced intake of macronutrients such as carbohydrates and protein is essential.\\
{\color{cFix}\textbf{C) The fly senses macronutrients through its taste system.}}\\
D) In the experiment, the researchers made sure the total calories in all three diets were the same.\\
E) The researchers tested the flies every two days for a week.\\
{\color{cFix}\textbf{F) Dopamine in the brain is closely related to diet.}}}
\par\smallskip
\textbf{正确答案}\ \textcolor{cFix}{\bfseries C, F}\qquad\textbf{你的答案}\ {\color{cWrong}\bfseries F}\qquad\textbf{本题得分：}\ {\color{cAccent}\bfseries 1 / 2}\quad{\footnotesize\color{gray}（用时 00:37；只选了 F，漏选 C）}
\end{qblock}

% ---------------- Q7 ----------------
\begin{qblock}{Q7\quad MCM-R \#114\ \ Persistent Back Pain}{多选 · 评分 \textcolor{cAccent}{\bfseries 0 / 2}}
\ptesub{文章 Passage}
\begin{promptbox}
People with persistent back pain or persistent headaches are twice as likely to suffer from both disorders, a new study from the University of Warwick has revealed.

The results, published in the Journal of Headache and Pain, suggest an association between the two types of pain that could point to a shared treatment for both.

The researchers from Warwick Medical School who are funded by the National Institute for Health Research (NIHR) led a systematic review of fourteen studies with a total of 460,195 participants that attempt to quantify the association between persistent headaches and persistent low back pain. They found an association between having persistent low back pain and having persistent (chronic) headaches, with patients experiencing one typically being twice as likely to experience the other compared to people without either headaches or back pain. The association is also stronger for people affected by migraine.

The researchers focused on people with chronic headache disorders, those who will have had headaches on most days for at least three months, and people with persistent low back pain that experience that pain day after day. These are two very common disorders that are leading causes of disability worldwide. Around one in five people have persistent low back pain and one in 30 have chronic headaches. The researchers estimate that just over one in 100 people (or well over half a million people) in the UK have both.
\end{promptbox}
\zh{华威大学一项新研究显示，患有持续性背痛或持续性头痛的人，同时患上这两种病症的可能性会翻倍。这项发表于《头痛与疼痛期刊》的研究结果，提示这两类疼痛之间存在关联，或可指向二者共同的治疗方法。由英国国家健康研究所（NIHR）资助的华威医学院研究人员，主导了一项系统综述，纳入14项研究、共46万零195名参与者，试图量化持续性头痛与持续性下背痛之间的关联。他们发现，患有持续性下背痛与患有持续性（慢性）头痛之间存在关联：与既无头痛也无背痛的人相比，患有其中一种病症的人患上另一种的可能性通常会翻倍；这种关联在偏头痛患者中更为明显。研究人员聚焦于患有慢性头痛病症的人（即在至少三个月里大多数日子都头痛的人），以及日复一日持续背痛的持续性下背痛患者。这是两种全球范围内导致残疾的主要常见病症：约五分之一的人患有持续性下背痛，三十分之一的人患有慢性头痛。研究人员估计，在英国，略超过百分之一的人（远超五十万人）同时患有这两种病症。}
\ptesub{题目 Question}
\begin{notebox}What are the main findings of the study?\end{notebox}
\ptesub{选项 Choices（正确项绿色加粗）}
{\small
{\color{cFix}\textbf{A) In the UK, one in 100 people have both persistent back pain and persistent headaches.}}\\
{\color{cFix}\textbf{B) There is a link between persistent back pain and persistent headache.}}\\
C) Headaches are the only cause of disability.\\
D) People with persistent back pain or headaches are half as likely to develop either condition.\\
E) The relationship between persistent headache and persistent low back pain is not quantifiable.}
\par\smallskip
\textbf{正确答案}\ \textcolor{cFix}{\bfseries A, B}\qquad\textbf{你的答案}\ {\color{cWrong}\bfseries E}\qquad\textbf{本题得分：}\ {\color{cAccent}\bfseries 0 / 2}\quad{\footnotesize\color{gray}（用时 00:24，全错）}
\end{qblock}

% ---------------- Q8 ----------------
\begin{qblock}{Q8\quad RO \#365\ \ Joint Venture}{段落重排 · 评分 \textcolor{cAccent}{\bfseries 1 / 3}}
\ptesub{乱序段落 Paragraphs}
\begin{promptbox}
\textbf{1)} A joint venture (JV) is a business arrangement in which two or more parties agree to pool their resources for the purpose of accomplishing a specific task.\\
\textbf{2)} In a JV, each of the participants is responsible for profits, losses, and costs associated with it.\\
\textbf{3)} However, the venture is its own entity, separate from the participants' other business interests.\\
\textbf{4)} This task can be a new project or any other business activity.
\end{promptbox}
\zh{（按正确顺序）合资企业（JV）是一种商业安排，其中两方或以上当事方同意集中各自的资源以完成一项特定任务；这项任务可以是一个新项目，也可以是其他任何商业活动。在合资企业中，每一位参与方都要对与之相关的利润、亏损和成本负责；不过，该合资企业本身是一个独立的实体，与参与方各自的其他商业利益相分离。}
\par\smallskip
\textbf{正确顺序}\ \textcolor{cFix}{\bfseries 1 → 4 → 2 → 3}\qquad\textbf{你的顺序}\ {\color{cWrong}\bfseries 1 → 2 → 3 → 4}\qquad\textbf{本题得分：}\ {\color{cAccent}\bfseries 1 / 3}\quad{\footnotesize\color{gray}（用时 00:35）}
\end{qblock}

% ---------------- Q9 ----------------
\begin{qblock}{Q9\quad RO \#369\ \ Mutations}{段落重排 · 评分 \textcolor{cAccent}{\bfseries 3 / 3}}
\ptesub{乱序段落 Paragraphs}
\begin{promptbox}
\textbf{1)} Others offer benefits, some can do both, and the mutation that underlies sickle cell disease is one that can be both good and very bad.\\
\textbf{2)} Our genes serve as an operating manual for cells of the body.\\
\textbf{3)} Genes tell cells what to do and when. But copying errors in those operating manuals - known as mutations - can lead to misspelled instructions that can change how cells operate.\\
\textbf{4)} Scientists now know that some of those mutations can lead to disease.
\end{promptbox}
\zh{（按正确顺序）我们的基因就像是身体细胞的操作手册。基因告诉细胞该做什么、何时做；但这些操作手册中的复制错误——即所谓的突变——可能导致指令出现"拼写错误"，从而改变细胞的运作方式。科学家现在知道，有些突变会导致疾病；另一些突变则带来益处，还有一些两者兼具——镰状细胞病背后的那个突变，就是一个既能带来好处、也能造成严重危害的突变。}
\par\smallskip
\textbf{正确顺序}\ \textcolor{cFix}{\bfseries 2 → 3 → 4 → 1}\qquad\textbf{你的顺序}\ \textcolor{cFix}{\bfseries 2 → 3 → 4 → 1}\qquad\textbf{本题得分：}\ {\color{cAccent}\bfseries 3 / 3}\quad{\footnotesize\color{gray}（用时 00:42，全对）}
\end{qblock}

% ---------------- Q10 ----------------
\begin{qblock}{Q10\quad FIBD\&D \#554\ \ Vegetative Propagation}{读填空（拖词）· 评分 \textcolor{cAccent}{\bfseries 0 / 4}}
\ptesub{文章 Passage（含填空作答）}
\begin{promptbox}
Because vegetative propagation is a form of asexual reproduction, plants produced through this system are genetic clones of a parent plant. One advantage of vegetative propagation is that plants with \rno{artificial}{favorable} traits are repeatedly reproduced. Commercial crop growers can employ \rno{traditional}{artificial} vegetative propagation techniques to ensure advantageous \rno{capabilities}{qualities} in their crops. A major disadvantage, however, of vegetative propagation is that it does not allow for any degree of genetic \rno{qualities}{variation}.
\end{promptbox}
\zh{由于营养繁殖是一种无性繁殖方式，通过这种方式产生的植物是母株的基因克隆体。营养繁殖的一个优点是，具有优良性状的植物可以被反复复制。商业作物种植者可以采用人工营养繁殖技术，来确保作物具备优良的性状。然而，营养繁殖的一个主要缺点是它完全不允许任何程度的基因变异。}
\ptesub{词库 Word Bank（拖词入空，正确答案加粗）}
{\small 1) \textbf{variation}\quad 2) \textbf{favorable}\quad 3) \textbf{artificial}\quad 4) capabilities\quad 5) diversification\quad 6) \textbf{qualities}\quad 7) traditional}
\par\smallskip
\textbf{你的答案：}\ {\color{cWrong}artificial, traditional, capabilities, qualities}\qquad\textbf{本题得分：}\ {\color{cAccent}\bfseries 0 / 4}\quad{\footnotesize\color{gray}（用时 02:18；四空全错，正确为 favorable / artificial / qualities / variation）}
\end{qblock}

% ---------------- Q11 ----------------
\begin{qblock}{Q11\quad FIBD\&D \#555\ \ Push-pull Factors}{读填空（拖词）· 评分 \textcolor{cAccent}{\bfseries 3 / 4}}
\ptesub{文章 Passage（含填空作答）}
\begin{promptbox}
In geographical terms, the push-pull factors drive people away from a place and \rok{draw} people to a new location. They help \rno{drag}{determine} migration of particular populations from one land to another. Push factors are often \rok{forceful}, demanding that a certain person or group of people leave one country for another. Pull factors, on the other hand, are often the positive policies that encourage people to \rok{immigrate} in order to seek a better life.
\end{promptbox}
\zh{在地理学意义上，推拉因素把人从一个地方推开、又把人吸引到一个新的地方。它们有助于决定特定人群从一片土地迁往另一片土地的迁移情况。推力因素往往带有强制性，迫使某个人或某群人离开一个国家前往另一个国家；而拉力因素则往往是那些鼓励人们为了寻求更好生活而移民的积极政策。}
\ptesub{词库 Word Bank（拖词入空，正确答案加粗）}
{\small 1) \textbf{immigrate}\quad 2) \textbf{forceful}\quad 3) drag\quad 4) \textbf{draw}\quad 5) \textbf{determine}\quad 6) formidable\quad 7) shift}
\par\smallskip
\textbf{你的答案：}\ draw, {\color{cWrong}drag}, forceful, immigrate\qquad\textbf{本题得分：}\ {\color{cAccent}\bfseries 3 / 4}\quad{\footnotesize\color{gray}（用时 01:43；第 2 空 drag→determine）}
\end{qblock}

% ---------------- Q12 ----------------
\begin{qblock}{Q12\quad FIBD\&D \#561\ \ India}{读填空（拖词）· 评分 \textcolor{cAccent}{\bfseries 1 / 4}}
\ptesub{文章 Passage（含填空作答）}
\begin{promptbox}
India's movie history began in 1896 when the Lumière brothers wanted to \rok{demonstrate} the art of cinema in Mumbai. Today, the films are known for their \rno{（未填）}{elaborate} singing and dancing. Indian dance, music and theater traditions \rno{drive}{span} back more than 2,000 years. The major classical dance traditions \rno{elaborate}{draw} on themes from mythology and literature and have rigid presentation rules.
\end{promptbox}
\zh{印度的电影历史始于1896年，卢米埃尔兄弟当时想在孟买展示电影这门艺术。如今，印度电影以其精心编排的歌舞著称。印度的舞蹈、音乐和戏剧传统可以追溯到2000多年前。主要的古典舞蹈传统取材于神话和文学主题，并有着严格的表演规范。}
\ptesub{词库 Word Bank（拖词入空，正确答案加粗）}
{\small 1) detailed\quad 2) \textbf{demonstrate}\quad 3) \textbf{draw}\quad 4) \textbf{span}\quad 5) drive\quad 6) \textbf{elaborate}\quad 7) appraise}
\par\smallskip
\textbf{你的答案：}\ demonstrate, {\color{cWrong}（第 2 空未作答）}, {\color{cWrong}drive, elaborate}\qquad\textbf{本题得分：}\ {\color{cAccent}\bfseries 1 / 4}\quad{\footnotesize\color{gray}（用时 02:20；第 2 空未填（应为 elaborate）、第 3 空 drive→span、第 4 空 elaborate→draw）}
\end{qblock}

% ---------------- Q13 ----------------
\begin{qblock}{Q13\quad FIBD\&D \#102\ \ Volcanoes}{读填空（拖词）· 评分 \textcolor{cAccent}{\bfseries 2 / 4}}
\ptesub{文章 Passage（含填空作答）}
\begin{promptbox}
Volcanoes blast more than 100 million tons of carbon dioxide into the atmosphere every year but the gas is usually \rok{harmless}. When a volcano erupts, carbon dioxide spreads out into the atmosphere and isn't \rok{concentrated} in one spot. But sometimes the gas gets trapped \rno{atmosphere}{underground} under enormous pressure. If it escapes to the surface in a dense \rno{underground}{cloud}, it can push out oxygen-rich air and become deadly.
\end{promptbox}
\zh{火山每年向大气层喷发超过一亿吨的二氧化碳，但这种气体通常是无害的。火山喷发时，二氧化碳会扩散到大气中，不会聚集在一个地点。但有时气体会在巨大的压力下被困在地下；如果它以浓密的云雾形式逃逸到地表，就会把富含氧气的空气挤走，从而变得致命。}
\ptesub{词库 Word Bank（拖词入空，正确答案加粗）}
{\small 1) \textbf{cloud}\quad 2) \textbf{concentrated}\quad 3) dangerous\quad 4) \textbf{harmless}\quad 5) \textbf{underground}\quad 6) air\quad 7) atmosphere\quad 8) collected\quad 9) over}
\ptesub{解析 Explanation（本套仅本题截图恰好展开了解析气泡）}
\begin{notebox}\footnotesize
\textbf{1.\ harmless}\quad 位于动词 is 后，选择形容词，可选的单词有 dangerous(危险的)、harmless(无害的)、harmful(有伤害的)；根据原文 "Volcanoes blast more than 100 million tons of carbon dioxide into the atmosphere every year but the gas is usually \_\_\_"，but 表示前后内容意思对立——前面说大量气体喷发，与之相反的是这些气体对我们没有伤害，故选 harmless。这里说的是：火山每年向大气喷发超过1亿吨的二氧化碳，但气体通常是无害的。【解题思路：主系表结构；转折关系；单句理解】\par\smallskip
\textbf{2.\ concentrated}\quad 位于 isn't 后，且主语为 carbon dioxide，可选择 v-ed 形式构成被动结构也可选择名词或形容词，可选的单词有 cloud(云)、concentrated(积聚)、dangerous(危险)、underground(在地下的)、aimed(瞄准)、air(空气)、harmful(有害的)、atmosphere(大气)、collected(收集)、fact(事实)。根据原文 "When a volcano erupts, carbon dioxide spreads out into the atmosphere and isn't \_\_\_ in one spot"，这里说二氧化碳进入大气层，在一个地点不会发生什么，代入选项，气体不会"聚集"在一个地方，最符合语意。【解题思路：被动语态；近义词辨析；单句理解】\par\smallskip
\textbf{3.\ underground}\quad 该空所在句主干："But sometimes the gas gets trapped \_\_\_ ..."，主干部分完整，因此选择副词作为修饰，可选择的单词有 underground(在地下)、over(倒下、翻转、穿过)。根据原文 "But sometimes the gas gets trapped \_\_\_ under enormous pressure"，由于二氧化碳在巨大的压力下，那应该是在地表下，而不是在水下承受巨大压力，因此选择 underground。这里说的是：但有时气体在巨大的压力下被困在地下。【解题思路：副词做状语；单句理解】\par\smallskip
\textbf{4.\ cloud}\quad 位于形容词 dense 后，选择名词，可选择的单词有 cloud(云)、air(空气)、atmosphere(大气)、collection(收集)、fact(事实)。根据原文 "If it escapes to the surface in a dense \_\_\_, it can push out oxygen-rich air and become deadly"，根据前面的冠词 a 判断，air 和 atmosphere 是不可数的，而 collection 又说不通，因此选择 cloud。这里说的是：如果它以浓稠云雾的形态逃逸出地表，就会排挤富含氧气的空气，从而变得致命。【解题思路：形容词+名词结构；单句理解】
\end{notebox}
\par\smallskip
\textbf{你的答案：}\ harmless, concentrated, {\color{cWrong}atmosphere, underground}\qquad\textbf{本题得分：}\ {\color{cAccent}\bfseries 2 / 4}\quad{\footnotesize\color{gray}（用时 03:18；第 3 空 atmosphere→underground、第 4 空 underground→cloud）}
\end{qblock}

% ---------------- Q14 ----------------
\begin{qblock}{Q14\quad FIBD\&D \#556\ \ Tiger Sharks}{读填空（拖词）· 评分 \textcolor{cAccent}{\bfseries 0 / 4}}
\ptesub{文章 Passage（含填空作答）}
\begin{promptbox}
The tiger shark is named for the dark, vertical stripes on either side of its body, which are reminiscent of a tiger's \rno{feature}{markings}. These stripes actually \rno{transplant}{fade} as the tiger shark ages, so they can't be used as an identifying \rno{markings}{feature} of every individual. Young tiger sharks have dark blotches or spots, which eventually \rno{fade}{merge} into stripes. For this reason, the species is sometimes known as the leopard shark or the spotted shark.
\end{promptbox}
\zh{虎鲨之所以得名，是因为它身体两侧有着让人联想到老虎斑纹的深色竖条纹。这些条纹实际上会随着虎鲨年龄的增长而褪色，因此不能作为辨识每一个个体的特征。幼年虎鲨身上有深色的斑块或斑点，这些斑点最终会融合成条纹。因此，这个物种有时也被称为豹纹鲨或斑点鲨。}
\ptesub{词库 Word Bank（拖词入空，正确答案加粗）}
{\small 1) transplant\quad 2) \textbf{merge}\quad 3) function\quad 4) \textbf{feature}\quad 5) \textbf{markings}\quad 6) \textbf{fade}\quad 7) detection}
\par\smallskip
\textbf{你的答案：}\ {\color{cWrong}feature, transplant, markings, fade}\qquad\textbf{本题得分：}\ {\color{cAccent}\bfseries 0 / 4}\quad{\footnotesize\color{gray}（用时 01:50；四空全错，正好错位一个：正确为 markings / fade / feature / merge）}
\end{qblock}

% ---------------- Q15 ----------------
\begin{qblock}{Q15\quad MCS-R \#160\ \ Walking Stamina}{单选 · 评分 \textcolor{cAccent}{\bfseries 1 / 1}}
\ptesub{文章 Passage}
\begin{promptbox}
If you've decided you want to improve your fitness, walking is a good choice. It's free, simple, and adaptable to your schedule. If you've been relatively sedentary, you might find that you can't walk very far at first without getting sore or out-of-breath. You just have to keep at it! If you try to walk a little further every day, you'll find that your walking stamina gradually improves. If you don't have the patience for that, there are a few other tricks you can try to help you reach your goals faster.
\end{promptbox}
\zh{如果你已经决定要改善自己的体能状况，步行是一个不错的选择。它免费、简单，还能灵活配合你的日程安排。如果你一直久坐少动，可能会发现一开始走不了多远就会酸痛或气喘。你只需要坚持下去！如果你每天都尝试多走一点，就会发现自己的步行耐力在逐渐提高。如果你没有耐心慢慢来，还有其他一些技巧可以帮你更快达成目标。}
\ptesub{题目 Question}
\begin{notebox}According to the passage, what might the ``tricks'' mentioned in the last sentence refer to?\end{notebox}
\ptesub{选项 Choices（正确项绿色加粗）}
{\small
{\color{cFix}\textbf{A) methods to improve your stamina quickly.}}\\
B) suggestions on how to keep fit.\\
C) benefits of different styles of walking.\\
D) different ways to build up your stamina over a long period.}
\par\smallskip
\textbf{正确答案}\ \textcolor{cFix}{\bfseries A}\qquad\textbf{你的答案}\ \textcolor{cFix}{\bfseries A}\qquad\textbf{本题得分：}\ {\color{cAccent}\bfseries 1 / 1}\quad{\footnotesize\color{gray}（用时 02:04，正确）}
\end{qblock}

% ---------------- Q16 ----------------
\begin{qblock}{Q16\quad MCS-R \#155\ \ Tree of Life}{单选 · 评分 \textcolor{cAccent}{\bfseries 1 / 1}}
\ptesub{文章 Passage}
\begin{promptbox}
Located in the waters of Homestead, Florida, Tree of Life is a serene villa complex designed to alleviate visitors from stress and anxiety. Iranian architect Milad Eshtiyaghi visualized the project as six independent villas connected by footpaths. Both the walkways and the buildings float gently on the existing pond and immerse guests in nature. To fit with the peaceful atmosphere of the surroundings, the villas are designed with organic forms, lacking harsh edges with the exception of their angled roofs. These organic buildings in Tree of Life were inspired by the vegetation in the surrounding landscape. So the complex could ``become one'' with nature in Homestead.
\end{promptbox}
\zh{坐落在佛罗里达州霍姆斯特德水域中的"生命之树"，是一个旨在缓解访客压力与焦虑的静谧别墅群。伊朗建筑师 Milad Eshtiyaghi 将该项目构想为六栋由步道相连的独立别墅。步道和建筑都轻盈地漂浮在现有的池塘上，让来客沉浸于自然之中。为了与周围宁静的氛围相契合，这些别墅采用了有机造型设计，除了斜屋顶外几乎没有锐利的棱角。"生命之树"中的这些有机建筑，其灵感来自周围景观中的植被，因此这一建筑群得以在霍姆斯特德与自然"融为一体"。}
\ptesub{题目 Question}
\begin{notebox}According to this passage, which one below is false?\end{notebox}
\ptesub{选项 Choices（正确项绿色加粗）}
{\small
A) Tree of Life is an ideal place for stressful people.\\
{\color{cFix}\textbf{B) All the buildings in this villa complex have harsh edges.}}\\
C) There are trees around the villa complex.\\
D) You will feel calm in the Tree of Life.}
\par\smallskip
\textbf{正确答案}\ \textcolor{cFix}{\bfseries B}\qquad\textbf{你的答案}\ \textcolor{cFix}{\bfseries B}\qquad\textbf{本题得分：}\ {\color{cAccent}\bfseries 1 / 1}\quad{\footnotesize\color{gray}（用时 01:35，正确）}
\end{qblock}

\end{document}
